\chapter{Proposta}\label{cap:proposta}

A proposta deste trabalho consistiu, essencialmente, na refatoração do sistema web E-Museu, desenvolvido inicialmente por \cite{gonzaga2024}, abrangendo aspectos de infraestrutura, código e testes, além da implementação de funcionalidades adicionais.

O desenvolvimento foi conduzido de forma metódica, com o uso de ferramentas do Github para o controle de tarefas \cite{github}, acompanhamento do fluxo de desenvolvimento, registro do histórico de modificações, bem como para o armazenamento e versionamento do código, contando com o apoio de professores nas questões necessárias.

Como forma de apresentação geral das modificações e implementações realizadas, destacam-se os seguintes itens:
\begin{itemize}
    \item Estabelecer uma estrutura de versionamento e desenvolvimento de código com padrões bem definidos, utilizando GitHub, de modo a garantir organização, rastreabilidade e colaboração eficiente.
    \item Configurar ambientes estáveis de desenvolvimento, teste e produção, assegurando consistência, reprodutibilidade e confiabilidade do sistema.
    \item Refatorar o código existente, aprimorando legibilidade, modularidade, manutenibilidade e desempenho. %, além de reestilizar a interface para proporcionar uma experiência mais clara, intuitiva e acessível aos usuários.
    \item Implementar testes automatizados a fim de assegurar a qualidade, confiabilidade e robustez do código, contemplando testes de fluxo, integração e usabilidade.
    %\item Desenvolver e aprimorar funcionalidades do E-Museu, como cadastro e gerenciamento de itens, upload de fotos e vídeos, geração de QR Code, relatórios de acesso e itens cadastrados, internacionalização do sistema e criação de newsletter, incorporando ainda mecanismos de segurança como reCAPTCHA e proteção de dados.
    \item Desenvolver e aprimorar funcionalidades do E-Museu, em particular cadastro com imagem de capa e galeria de fotos, internacionalização do catálogo, incorporação de mecanismo antibot baseado em desafio com \textit{token} (Cloudflare Turnstile) e geração de código QR para apoio à identificação e etiquetagem das peças físicas.
    \item Aplicar padrões de projeto e boas práticas de desenvolvimento, orientando tanto a refatoração quanto a implementação de novas funcionalidades de forma estruturada e eficiente.
\end{itemize}

O trabalho a ser desenvolvido deverá impactar positivamente para diferentes públicos. Os profissionais e alunos que atuam e futuramente atuarão na administração e catalogação de itens; a professora responsável pela administração do site; os desenvolvedores, que contarão com um código refatorado e testado; e, por fim, os usuários finais, que poderão acessar o E-Museu com novas funcionalidades.

Ao final do trabalho, entregou-se uma versão refatorada e aprimorada do sistema web do E-Museu, com infraestrutura organizada, código evoluído, testes automatizados e novas funcionalidades implementadas. Como benefícios, observa-se melhoria da qualidade e confiabilidade do sistema, ampliação de recursos do catálogo, fortalecimento da segurança em fluxos sensíveis e garantia de uma base técnica que possibilita manutenção e evolução contínuas do E-Museu.

